% Study sheet: Volume I, Chapter 7, ML Frameworks.
% Numbers: \val{...} from frameworks_studysheet_numbers.tex, written by frameworks_studysheet_numbers.py.
\documentclass{studysheet}
\input{frameworks_studysheet_numbers}
\sheetvolume{1}
\sheetchapter{7}{ML Frameworks}
\sheetsources{Numbers from the mlsysim registry and the chapter's scenarios; per-operation costs are the chapter's model; assumed inputs are marked.}

\begin{document}
% ===================================================================== FRONT
\frontside
\begin{bigidea}
A framework turns a model into \hl{kernels on a machine}. The model fixes the work; the framework
decides when each operation runs, what training keeps, and which kernel runs where. Those choices set
each step's \hl{overhead, bytes and efficiency}, often several-fold.
\end{bigidea}

\begin{leftcol}
\sheetsec{Key ideas}
\begin{ideas}
\idea{Three problems, three terms.}{Execution sets overhead $\Latency$; differentiation adds bytes
  $\DataVolume$ and work $\Operations$; abstraction (kernel, layout, precision) sets efficiency $\EfficiencyHw$.}
\idea{A module carries a mode.}{Evaluation mode turns off dropout and freezes batch-norm statistics;
  no-grad stops saving activations. Validation needs both.}
\idea{Every operation costs host time.}{Python dispatch plus a kernel launch, microseconds each.
  Shorter kernels leave the accelerator waiting.}
\idea{Fusion cuts trips to memory.}{A fused chain reads and writes its tensor once, not once per
  operation, and launches once; it still stores what backward reads.}
\idea{Gradients cost memory.}{One reverse sweep yields every weight gradient; its saved activations grow
  with batch and depth; micro-batches or checkpointing cut the peak.}
\idea{The host link is slow.}{Pinned memory on a second stream overlaps copies with compute;
  one shared stream, or a host that waits on each copy, serializes them.}
\end{ideas}

\vspace{10pt}
\sheetsec{One eager decoding step, to scale}
{\centering
\begin{tikzpicture}[sheet]
  \pgfmathsetmacro\u{21}
  \pgfmathsetmacro\xp{56+\val{dpy}*\u}
  \pgfmathsetmacro\xh{56+\val{dhost}*\u}
  \pgfmathsetmacro\xr{56+\val{dread}*\u}
  \node[lab, anchor=east] at (52,52) {host issues};
  \draw[accbox] (56,40) rectangle (\xp,64);
  \draw[hotbox] (\xp,40) rectangle (\xh,64);
  \node[lab, anchor=center] at ({(56+\xp)/2},52) {Python \val{dpy} ms};
  \node[lab, text=white, anchor=center] at ({(\xp+\xh)/2},52) {launch \val{dln} ms};
  \node[note, anchor=south east] at (\xh,68) {\val{dkern} kernels $\times$ (\val{pyus} $+$ \val{lnus}) $\mu$s $=$ \val{dhost} ms};
  \node[lab, anchor=east] at (52,12) {device reads};
  \draw[nbox] (56,0) rectangle (\xr,24);
  \node[note, anchor=west] at (\xr+4,12) {\val{dgb} GB of weights at \val{dbw} TB/s: \val{dread} ms};
  \draw[dim] (56,-4) -- (56,30);
\end{tikzpicture}\par}
\vspace{3pt}
{\centering\fontsize{8.6}{11}\selectfont GPT-2 XL, batch one, \val{dlayers} blocks $\times$ \val{dperblock} kernels.
Replaying captured launches as one submission leaves about the weight read.\par}
\end{leftcol}%
\begin{rightcol}
\sheetsec{Counting rules}
\fcard{Dispatch ratio}%
  {\text{ratio} = \frac{\HostSoftwareTime}{\IronLawComputeTime + \IronLawMemoryTime}}%
  {Host time $\HostSoftwareTime$ $=$ Python calls $\times$ \val{pyus} $\mu$s $+$ kernel launches $\times$ \val{lnus} $\mu$s
   (eager: one of each per operation). Ideal: $\IronLawComputeTime = \Operations/\PeakRate$,
   $\IronLawMemoryTime = \DataVolume/\Bandwidth$.}%
  {Ratio $>1$: overhead-bound, dispatch share ratio$/(1+$ratio$)$. Otherwise the larger device term rules.}{}
\vspace{8pt}
\begin{sheettable}{@{}l L@{}}
\thead{Cost} & \thead{Rule}\\ \midrule
Fused chain & $2\ChainedOperationCount\TensorElementCount\ElementStorageSize$ bytes unfused, $2\TensorElementCount\ElementStorageSize$ fused\\
Compile breakeven & $\CompilationBreakevenCount = \dfrac{\CompilationTime}{\EagerExecutionTime - \CompiledExecutionTime}$ runs\\
Saved activations & $\ActivationMemory \approx \BatchSize\,\LayerCount\,\SequenceLength\,\HiddenDimension\,\LayerSavedTensorCount\,\ElementStorageSize$\\
Host copy & $\BatchTransferTime \approx$ batch bytes$/\HostLinkBandwidth$ (startup negligible); per batch, serial compute $+ \BatchTransferTime$, overlapped $\max($compute$, \BatchTransferTime)$\\
\end{sheettable}
\vspace{3pt}
{\fontsize{8.4}{10.4}\selectfont\color{muted}$\ChainedOperationCount$ elementwise operations on $\TensorElementCount$ values of $\ElementStorageSize$ bytes;
times per run $\EagerExecutionTime$, $\CompiledExecutionTime$; batch $\BatchSize$, layers $\LayerCount$,
sequence $\SequenceLength$, width $\HiddenDimension$, $\LayerSavedTensorCount$ tensors saved per layer; $\HostLinkBandwidth$ host-link bandwidth.\par}

\vspace{12pt}
\sheetsec{Choosing a path}
\begin{sheettable}{@{}L L@{}}
\thead{If} & \thead{Then}\\ \midrule
code changes every few minutes & stay eager\\
the same graph runs more than $\CompilationBreakevenCount$ times & compile (ahead of time for a fixed target)\\
ratio $> 1$ & fewer operations (fuse), cheaper launches (replay a captured graph), or a bigger batch (same host cost, more device work)\\
choosing a framework & drop paths that cannot export every operator to a runtime that fits the target; then weigh team pipeline and migration cost\\
\end{sheettable}
\end{rightcol}

% ===================================================================== BACK
\backside
\begin{answers}
\ans{Finish}{\val{yentries} $\times$ \val{pyus} $+$ \val{ylaunch} $\times$ \val{lnus} $=$ \val{ythost} $\mu$s;
  ratio $\val{ythost} \div \val{wtdev} \approx \val{yratio}$, still overhead-bound (\val{yshare}\% dispatch).
  After replay, memory (\val{wtm} $\mu$s) $>$ compute (\val{wtc} $\mu$s): memory-bound next (dispatch card).}
\ans{1}{Unfused $\val{q1rw} \times \val{q1k} \times \val{q1n} \times \val{q1bytes}$ bytes $=$ \val{q1unf} MB, \val{q1tunf} $\mu$s; fused
  \val{q1fus} MB, \val{q1tfus} $\mu$s. Two fewer operations save \val{q1saved} $\mu$s of host time.}
\ans{2}{\val{q2comp} s $\div$ (\val{q2te} $-$ \val{q2tc}) ms $=$ \val{q2be} runs. Serving passes it in
  \val{q2min} min: compile. A notebook's \val{q2nb} runs never do: stay eager (breakeven).}
\columnbreak
\ans{3}{\val{q3b} $\times$ \val{q3l} $\times$ \val{q3s} $\times$ \val{q3d} $\times$ \val{q3k} $\times$ \val{q3bytes} bytes
  $=$ \val{q3gb} GB; batch \val{q3b2}: \val{q3gb2} GB $>$ \val{q3free} GB, no. Accumulate gradients over two
  micro-batches of \val{q3b}, or checkpoint (saved activations).}
\ans{4}{\val{q4mb} MB $\div$ \val{q4link} GB/s $=$ \val{q4t} ms. Serial \val{q4t} $+$ \val{q4c} $=$ \val{q4serial} ms;
  overlapped \val{q4ovl} ms, set by the copy. No overlap: one shared stream, or a
  host that waits on each copy (host-link idea).}
\end{answers}

\begin{leftcol}
\sheetsec{Worked example}
\begin{practicebox}\boxtitle{Is a small training step limited by dispatch?}
The chapter's two-layer MLP (\val{win}$\,\to\,$\val{whid}$\,\to\,$\val{wout}, batch \val{wbatch}, 32-bit):
\val{wmflop} MFLOP and \val{wmb} MB per step, on \val{wpeak} TFLOP/s and \val{wbw} TB/s. Eager:
\val{wentries} Python entries (\val{wfwd} forward operations, one backward call), \val{wlaunch} launches.
\begin{steps}
\item Device: \val{wmflop} MFLOP $\div$ \val{wpeak} TFLOP/s $=$ \val{wtc} $\mu$s; \val{wmb} MB $\div$ \val{wbw} TB/s
  $=$ \val{wtm} $\mu$s; together \res{\val{wtdev} $\mu$s} (ideal).
\item Host: \val{wentries} calls $\times$ \val{pyus} $+$ \val{wlaunch} $\times$ \val{lnus} $=$ \res{\val{wthost} $\mu$s}.
\item Ratio: \val{wthost} $\div$ \val{wtdev} $\approx$ \res{\val{wratio}}; dispatch share \val{wshare}\%.
\end{steps}
\takeaway{Count launches before FLOPs: at this size the host sets the time.}
\end{practicebox}

\vspace{12pt}
\sheetsec{Your turn: finish it}
\begin{practicebox}\boxtitle{Does compiling the step fix it?}
\assumed{} compiling fuses the step into \val{ylaunch} kernels behind \val{yentries} Python call;
device time stays \val{wtdev} $\mu$s.\par\vspace{4pt}
\textbf{Given:} host time $=$ \val{yentries} $\times$ \val{pyus} $+$ \val{ylaunch} $\times$ \val{lnus} $\mu$s.\par
\textbf{You:} find the ratio and the dispatch share. Still overhead-bound? If replay then made host time
negligible, which device term would set the step?\par\vspace{50pt}
\end{practicebox}
\end{leftcol}%
\begin{rightcol}
\sheetsec{Common mistakes}
\mistake{``Same model, same math, so every framework runs it at the same speed.''}%
  {The eager decoding step spends \val{dhost} ms issuing kernels; replayed as one submission, it
   falls to about the \val{dread} ms weight read, \val{dratio}$\times$ less time.}{}
\mistake{``Compilation is a one-time cost, so always turn it on.''}%
  {Illustrative ResNet-50: \val{m2img} images with \val{m2rec} recompilations take
   \val{m2eager} s eager and \val{m2compiled} s compiled (\val{m2comp} s per compile): \val{m2slow}$\times$ slower.}{}
\mistake{``I wrapped evaluation in no-grad, so the model is in evaluation mode.''}%
  {No-grad only stops saving for backward. Dropout still drops until the module is set to evaluation
   mode, so predictions on fixed inputs vary run to run. The reverse slip, evaluation mode
   without no-grad, saves activations: peak memory rises.}{}

\vspace{2pt}
\sheetsec{Check yourself}
\begin{qs}
\item \assumed{} eager bias add, dropout and ReLU on \val{q1n} 16-bit values at \val{wbw} TB/s. Bytes and time,
  unfused and fused? What else does fusion save?
\item \assumed{} a model serves in \val{q2te} ms eager and \val{q2tc} ms compiled; compiling takes \val{q2comp} s.
  Compile for serving at \val{q2rps} requests/s? For a notebook that runs \val{q2nb} times between edits?
\item \assumed{} training with $\BatchSize =$ \val{q3b}, $\LayerCount =$ \val{q3l}, $\SequenceLength =$ \val{q3s},
  $\HiddenDimension =$ \val{q3d}, \val{q3k} saved tensors per layer, 16-bit. Saved activations? Does batch
  \val{q3b2} fit in \val{q3free} GB free for activations? If not, two ways out.
\item \assumed{} each step copies a \val{q4batch}-image 32-bit batch (\val{q4mb} MB) over a
  \val{q4link} GB/s host link and computes for \val{q4c} ms. Time per batch after warm-up (no startup), serial and
  overlapped? A timeline shows copies and kernels never running at once: two causes?
\end{qs}
\end{rightcol}
\end{document}
